\daytitle{4}{运动图像与追及相遇}

\begin{learningbox}{本节导读}
$x$-$t$ 图和 $v$-$t$ 图是运动学的“可视化语言”。本节学习读图、画图，并把追及相遇中的位置关系、速度关系和距离变化放到图像中判断。
\end{learningbox}

\textbf{本节读题路线：} 先分清图像是 $x$-$t$ 还是 $v$-$t$；再读斜率、面积、截距和交点各代表什么。追及相遇题把两物体画在同一坐标图中：位置相等决定是否相遇；速度相等只是距离可能取极值的候选条件，还要结合速度关系的变化和考察区间判断。

\section*{一、学习目标}

\begin{enumerate}
  \item 能读懂 $x$-$t$ 图像，知道斜率、截距、折点和曲线切线的物理意义。
  \item 能读懂 $v$-$t$ 图像，会用有向面积求区间位移 $\Delta x$、用斜率求加速度。
  \item 会区分匀速、匀加速、匀减速、静止四种运动在不同图像中的形态。
  \item 能把追及相遇问题中的位置关系和速度关系同时放进图像与方程中分析。
  \item 会判断追及过程中的距离极值（最远或最近）。
\end{enumerate}

\section*{二、初高中衔接}

初中物理讲过“匀速直线运动的路程-时间图像”，常见的是一条过原点的直线。高中要在这个基础上升级三层读图能力：

\begin{enumerate}
  \item 从“路程随时间变化”升级到“位置随时间变化”，能从斜率读出速度的大小和方向。
  \item 增加 $v$-$t$ 图：斜率表示加速度，图线与时间轴围成的有向面积对应区间位移 $\Delta x$。
  \item 把图像用于分析过程：找速度相等、停车、相遇等关键时刻，再判断距离怎样变化。
\end{enumerate}

\textbf{科学发展视角：} 真实实验和工程系统常常先得到一串数据点，再用图像寻找规律。车辆测速、运动员训练分析、物流轨迹监控和机器人路径规划，都要从位置-时间、速度-时间数据中读出斜率、面积和关键时刻。图像法不是“辅助画图”，而是把数据变成判断的工具。

\section*{三、核心概念}

\begin{conceptbox}{$x$-$t$ 图像（位置-时间图）}

\textbf{直观理解：} $x$-$t$ 图像描述的是物体“位置随时间怎么变”。

\begin{center}
% TikZ 图：x-t 图像示例（匀速直线）
% Day 04：运动图像与追及相遇
\begin{tikzpicture}[scale=1,>=Latex]
  \draw[->] (-0.3,0) -- (5.5,0) node[right] {$t$};
  \draw[->] (0,-0.3) -- (0,4.5) node[above] {$x$};
  \node[below left] at (0,0) {$O$};
  % 匀速直线：x = 0.8t
  \draw[very thick] (0,0) -- (5,4);
  \node[above left] at (2.5,2) {$x = vt$};
  % 标注斜率
  \draw[dashed] (3,0) -- (3,2.4);
  \draw[dashed,thick] (0,0) -- (3,2.4);
  \node[below] at (1.5,0) {$t$};
  \node[left] at (0,1.2) {$x$};
  \node[above right] at (0.5,0.8) {斜率 $= v$};
\end{tikzpicture}


\figcaption{匀速直线运动的 $x$-$t$ 图像：位置随时间均匀变化，斜率 $=v$}\label{fig:xt-graph}
\end{center}

\textbf{关键读图方法：}
\begin{itemize}
  \item \textbf{斜率} = 速度 $v$。斜率的绝对值越大，速率越大；斜率为负表示沿所选正方向的反方向运动。
  \item \textbf{截距} = 初始位置 $x_0$。图线与纵轴的交点。
  \item \textbf{交点} = 两物体在该时刻位置相同（相遇）。
  \item \textbf{水平线} = 物体静止。
  \item \textbf{非直线曲线} = 变速运动（各点切线的斜率为相应时刻的瞬时速度）。
\end{itemize}

\textbf{实验数据视角：} 真实实验中，位置和时间往往是一组离散数据点，例如打点计时器、光电门或手机传感器记录的数据。图像中的直线或曲线，是根据数据建立的运动模型。读斜率、读面积，就是从数据中提取速度、加速度和位移信息。

\end{conceptbox}

\begin{conceptbox}{$v$-$t$ 图像（速度-时间图）}

\textbf{直观理解：} $v$-$t$ 图像描述的是“速度随时间怎么变”。

\textbf{关键读图方法（回顾 Day 03）：}
\begin{itemize}
  \item \textbf{斜率} = 加速度 $a$。
  \item \textbf{有向面积} 对应所选时间区间内的位移 $\Delta x$（$t$ 轴上方面积为正，下方面积为负）。
  \item \textbf{纵坐标} = 瞬时速度。
  \item \textbf{横线} = 匀速运动。
  \item \textbf{倾斜直线} = 匀变速运动。
  \item \textbf{交点} = 两物体在该时刻速度相等（不一定相遇）。
\end{itemize}

\begin{center}
% TikZ 图：常见运动的 v-t 图像对比
% Day 04：运动图像与追及相遇
% 分成三个小坐标系，避免三条线交于一点造成误读。
\begin{tikzpicture}[scale=1,>=Latex]
  \foreach \x/\title in {0/匀速,2.5/匀加速,5/匀减速} {
    \begin{scope}[xshift=\x cm]
      \draw[->] (0,0) -- (1.8,0) node[right] {$t$};
      \draw[->] (0,0) -- (0,1.8) node[above] {$v$};
      \node[below left] at (0,0) {$O$};
      \node[below] at (0.9,-0.35) {\title};
    \end{scope}
  }

  \draw[very thick] (0.25,1.0) -- (1.55,1.0);
  \begin{scope}[xshift=2.5cm]
    \draw[very thick] (0.25,0.35) -- (1.55,1.45);
  \end{scope}
  \begin{scope}[xshift=5cm]
    \draw[very thick] (0.25,1.45) -- (1.55,0.35);
  \end{scope}
\end{tikzpicture}


\figcaption{三种典型运动的 $v$-$t$ 图像对比}\label{fig:vt-comparison}
\end{center}

\end{conceptbox}

\section*{四、公式与适用条件}

\textbf{读图前检查：}
\begin{enumerate}
  \item 先看横轴、纵轴和单位，不要把 $x$-$t$ 图误读成 $v$-$t$ 图。
  \item 斜率、面积、交点、截距的物理意义只由图像类型决定。
  \item 两物体比较时，必须确认它们是否在同一参考系和同一时间轴上。
  \item 追及题先写位置关系，再用速度相等判断距离极值。
\end{enumerate}

\begin{learningbox}{读运动图像的四步分析清单}
看到图像时，不要先问“这条线像什么”，而要按顺序读信息：

\begin{center}
\begin{tabular}{lll}
\toprule
步骤 & 先看什么 & 得到什么物理信息 \\
\midrule
1 & 横轴、纵轴和单位 & 判断是 $x$-$t$ 图还是 $v$-$t$ 图 \\
2 & 起点、终点、截距 & 初位置、初速度或初始条件 \\
3 & 斜率的正负和大小 & 速度或加速度的方向与大小 \\
4 & 交点、面积、折点或曲线形状 & 相遇、位移、运动状态变化 \\
\bottomrule
\end{tabular}
\end{center}

\textbf{关键提醒：} 同样是一条直线，在 $x$-$t$ 图中表示匀速运动；在 $v$-$t$ 图中表示匀变速运动。同样是交点，在 $x$-$t$ 图中表示位置相同；在 $v$-$t$ 图中只表示速度相同。
\end{learningbox}

\subsection*{两种图像的核心对比}

\begin{center}
\begin{tabular}{lll}
\toprule
 & $x$-$t$ 图 & $v$-$t$ 图 \\
\midrule
纵轴 & 位置 $x$ & 速度 $v$ \\
斜率 & 速度 $v$ & 加速度 $a$ \\
面积 & 高中运动学中不直接表示常用物理量 & 有向面积对应位移 $\Delta x$（上正下负） \\
水平线 & 静止 & 匀速 \\
交点 & 相遇 & 速度相等 \\
截距 & 初始位置 & 初速度 \\
\bottomrule
\end{tabular}
\end{center}

\subsection*{追及问题的核心关系}

追及题先把初始位置说清楚。若追者在后、前车领先 $\Delta x$，并且两车位移都从各自初始位置开始计，则：
\begin{itemize}
  \item \textbf{追上条件：} 追者多走的位移等于初始距离，$x_{\text{追}} - x_{\text{前}} = \Delta x$。
  \item \textbf{相距最远/最近的内部候选条件：} 若距离的极值出现在所研究时间区间内部，通常先找 $v_{\text{追}} = v_{\text{前}}$；随后检查相对速度是否变号。初始、停车、相遇和考察终点也必须比较。
  \item \textbf{$v$-$t$ 图像法：} 按统一的“前减后”或“后减前”顺序取两条速度图线之差，其有向面积就是相对位移的变化量。
\end{itemize}

\textbf{为什么速度相等是距离极值的候选条件？}

设前车与后车的距离为 $s=x_{\text{前}}-x_{\text{后}}$。把速度理解为某一时刻的瞬时速度时，距离的瞬时变化率\textbf{严格等于} $v_{\text{前}}-v_{\text{后}}$：前车更快时，$s$ 增大；后车更快时，$s$ 减小。当 $v_{\text{前}}=v_{\text{后}}$ 时，距离暂时不增不减，这是距离取极值的\textbf{候选时刻}。还必须检查速度关系是否确实由“前车更快”变成“后车更快”，或反向变化；只有发生这种变号，才分别得到最大距离或最小距离。若相对速度不变号，极值可能出现在初始、停车、相遇或考察终点，而不是速度相等时刻。

\begin{center}
% TikZ 图：两物体追及的 x-t 图
\begin{tikzpicture}[scale=0.9,>=Latex,line cap=round,line join=round]
  \draw[->] (-0.25,0) -- (6.15,0) node[right] {$t$};
  \draw[->] (0,-0.25) -- (0,5.05) node[above] {$x$};
  \node[below left] at (0,0) {$O$};

  % 前车初始位置较大、速度较小；后车初始位置较小、速度较大。
  \draw[very thick,red!75!black] (0,1.6) -- (5.6,3.84)
    node[right] {前车};
  \draw[very thick,blue!75!black] (0,0.4) -- (5.6,4.88)
    node[right] {后车};

  % t_1 时刻，两条图线的纵向间隔就是两车距离。
  \coordinate (RearOne) at (1.6,1.68);
  \coordinate (FrontOne) at (1.6,2.24);
  \draw[dashed,gray!80] (1.6,0) -- (FrontOne);
  \draw[<->,thick] (RearOne) -- (FrontOne)
    node[midway,left,xshift=-2pt] {$s(t_1)$};
  \node[below] at (1.6,0) {$t_1$};

  % 两图线交点表示位置相同。
  \coordinate (Meet) at (3,2.8);
  \fill (Meet) circle (1.8pt);
  \draw[dashed,gray!80] (3,0) -- (Meet);
  \node[below] at (3,0) {$t_{\text{相}}$};
  \node[above left] at (Meet) {$P$};
  \node[align=left,anchor=west] at (3.28,2.62)
    {\small 交点：$x_{\text{前}}=x_{\text{后}}$\\[-1pt]\small 两车相遇};
\end{tikzpicture}


\figcaption{两物体的 $x$-$t$ 图：同一时刻的纵向间隔表示距离，图线交点表示相遇}\label{fig:xt-chase-two-objects}
\end{center}

\begin{center}
% TikZ 图：追及过程中距离增减与相对速度的关系
\begin{tikzpicture}[scale=0.78,>=Latex,line cap=round,line join=round,
  every node/.style={font=\small}]
  % 三幅图表示依次发生的三个时刻；两车都向右运动。
  \foreach \x in {0,5.1,10.2} {
    \draw[gray!65] (\x,0) -- ++(4.2,0);
  }

  % 极值前：前车更快，距离增大。
  \filldraw[fill=blue!12,draw=blue!70!black] (0.25,0.02) rectangle +(0.72,0.34);
  \filldraw[fill=red!12,draw=red!70!black] (2.45,0.02) rectangle +(0.72,0.34);
  \draw[->,blue!75!black,thick] (0.61,0.61) -- +(0.72,0)
    node[above] {$v_{\text{后}}$};
  \draw[->,red!75!black,thick] (2.81,0.61) -- +(1.12,0)
    node[above] {$v_{\text{前}}$};
  \draw[<->] (0.97,0.18) -- (2.45,0.18) node[midway,above] {$s_1$};
  \node[align=center] at (2.1,-0.74) {极值前：$v_{\text{后}}<v_{\text{前}}$\\距离增大};

  % 候选时刻：速度相等，相对速度为零。
  \filldraw[fill=blue!12,draw=blue!70!black] (5.65,0.02) rectangle +(0.72,0.34);
  \filldraw[fill=red!12,draw=red!70!black] (8.15,0.02) rectangle +(0.72,0.34);
  \draw[->,blue!75!black,thick] (6.01,0.61) -- +(0.94,0)
    node[above] {$v_{\text{后}}$};
  \draw[->,red!75!black,thick] (8.51,0.61) -- +(0.94,0)
    node[above] {$v_{\text{前}}$};
  \draw[<->] (6.37,0.18) -- (8.15,0.18) node[midway,above] {$s_*$};
  \node[align=center] at (7.2,-0.74) {候选时刻：$v_{\text{后}}=v_{\text{前}}$\\需检查相对速度是否变号};

  % 极值后：后车更快，距离减小。
  \filldraw[fill=blue!12,draw=blue!70!black] (11.25,0.02) rectangle +(0.72,0.34);
  \filldraw[fill=red!12,draw=red!70!black] (13.20,0.02) rectangle +(0.72,0.34);
  \draw[->,blue!75!black,thick] (11.61,0.61) -- +(1.18,0)
    node[above] {$v_{\text{后}}$};
  \draw[->,red!75!black,thick] (13.56,0.61) -- +(0.72,0)
    node[above] {$v_{\text{前}}$};
  \draw[<->] (11.97,0.18) -- (13.20,0.18) node[midway,above] {$s_2$};
  \node[align=center] at (12.3,-0.74) {极值后：$v_{\text{后}}>v_{\text{前}}$\\距离减小};
\end{tikzpicture}


\figcaption{追及过程的距离变化：相对速度决定距离增减，速度相等只是极值候选时刻}\label{fig:chase-scenarios}
\end{center}

\section*{五、典型模型}

追及问题应先画清四个节点：初始距离、速度相等时刻、可能停车时刻、追上时刻。速度相等是区间内部距离极值的候选时刻；是否真的取到极值、是否追上，还要检查相对速度是否变号并回到位置关系判断。

\begin{center}
\renewcommand{\arraystretch}{1.18}
\begin{tabular}{p{0.18\textwidth}p{0.32\textwidth}p{0.36\textwidth}}
\toprule
模型 & 关键判断 & 方程入口 \\
\midrule
匀速追匀速 & 后车速度更大才可能追上 & $(v_A-v_B)t=\Delta x$ \\
匀加速追匀速 & 后车初速度 $v_0<v_B$ 时，初期距离增大；相对速度由负变正时距离最大 & 追上：$v_0t+\frac12 at^2=v_Bt+\Delta x$；$v_0=0$ 时再化简 \\
匀速追匀减速 & 若前车按题意减速至停下并保持静止，后车持续前进则最终一定追上；先判是否在停车前追上 & $t_0=v_B/a$；停车前后分段比较相对位移 \\
\bottomrule
\end{tabular}
\end{center}

\noindent 使用顺序：先写位置关系，再找速度相等或停车这样的关键时刻，最后比较相对位移与初始距离。这样比孤立记公式更可靠。

\subsection*{追及题列式框架}

\begin{enumerate}
  \item \textbf{定正方向和初始位置：} 通常取运动方向为正，把领先距离写成 $\Delta x$ 或 $d$，不要默认两车从同一点出发。
  \item \textbf{分别写位移：} 用 $x_A(t)$、$x_B(t)$ 表示两车从各自初始位置起算的位移。
  \item \textbf{写追上条件：} 后车追上前车时，后车多走的位移等于初始距离，即 $x_{\text{后}}-x_{\text{前}}=\Delta x$。
  \item \textbf{找极值候选：} 先找速度相等、停车、相遇和考察终点；若速度相等时相对速度发生变号，它才对应区间内部的最大或最小距离。
\end{enumerate}

\textbf{使用提醒：} 这个框架用来防止漏过程，不表示每道题都要写满四步。若题目出现停车、反向运动或分段速度图像，应优先把过程分清；若速度相等时刻不在实际过程内，或相对速度没有变号，就不能把它直接当作距离极值。

\section*{六、例题精讲}

\begin{examplebox}{读 $x$-$t$ 图}

\textbf{题意分析：} 某物体沿直线运动的 $x$-$t$ 图像如图。

\begin{center}
% TikZ 图：x-t 折线图（OA匀速→AB静止→BC返回）
% Day 04：运动图像与追及相遇
% A、B、C 标在折线拐点和终点
\begin{tikzpicture}[scale=1,>=Latex]
  \draw[->] (-0.3,0) -- (5.5,0) node[right] {$t$};
  \draw[->] (0,-0.3) -- (0,4) node[above] {$x$};
  \node[below left] at (0,0) {$O$};
  % O点
  \filldraw[black] (0,0) circle (2pt);
  % OA: 匀速（标A在拐点）
  \draw[very thick] (0,0) -- (2,3);
  \filldraw[black] (2,3) circle (2pt) node[above right] {$A$};
  % AB: 静止（标B在拐点）
  \draw[very thick] (2,3) -- (3.5,3);
  \filldraw[black] (3.5,3) circle (2pt) node[above] {$B$};
  % BC: 返回（标C在终点）
  \draw[very thick] (3.5,3) -- (5,0);
  \filldraw[black] (5,0) circle (2pt) node[above right] {$C$};
  % 标注
  \node[below] at (2,0) {$t_A$};
  \node[below] at (3.5,0) {$t_B$};
  \node[below] at (5,0) {$t_C$};
  \node[left] at (0,3) {$x_A$};
\end{tikzpicture}


\figcaption{某物体的 $x$-$t$ 图像：$OA$ 匀速 $\to$ $AB$ 静止 $\to$ $BC$ 返回}\label{fig:xt-brokenline}
\end{center}

$OA$ 段为过原点的直线，$AB$ 段为水平线，$BC$ 段为倾斜直线（斜率为负）。分析各段的运动情况。

\textbf{物理模型：} $x$-$t$ 图分析。

\textbf{解题过程：}
\begin{itemize}
  \item $OA$ 段：物体从原点出发做匀速直线运动，速度 $v = \dfrac{x_A}{t_A}$。
  \item $AB$ 段：物体静止（$x$ 不变），速度为零。
  \item $BC$ 段：物体沿原路返回，有符号速度
  \[
  v_{BC}=\frac{x_C-x_B}{t_C-t_B}<0,
  \]
  速率为 $\lvert v_{BC}\rvert=\dfrac{|x_C-x_B|}{t_C-t_B}$。
  \item $C$ 点：物体回到初始位置，因此从 $O$ 到 $C$ 的总位移为零。
\end{itemize}

\textbf{方法总结：} 读 $x$-$t$ 图的要点：先看斜率（正负和绝对值），再看截距（起点位置），最后看折点或曲线形状怎样变化。

\end{examplebox}

\begin{examplebox}{追及问题——公式法}

\textbf{题意分析：} 一辆汽车（$A$）以 $v_A = 20\ \mathrm{m/s}$ 行驶，发现前方 $d = 60\ \mathrm{m}$ 处有一辆摩托车（$B$）以 $v_B = 10\ \mathrm{m/s}$ 同向匀速行驶。汽车立即以大小为 $a = 2\ \mathrm{m/s^2}$ 的加速度匀减速。问两车是否相撞？

\textbf{物理模型：} 匀减速（汽车）追匀速（摩托车）。这里的 $a$ 表示减速度大小，判断是否追尾：当汽车速度减到与摩托车速度相等时，比较汽车位移与摩托车位移 $+ d$。

\textbf{解题过程：}

速度相等时：
\[
v_A - a t = v_B,\qquad 20 - 2t = 10,\qquad t = 5\ \mathrm{s}
\]

此时汽车位移：
\[
x_A = v_A t - \frac12 a t^2 = 20 \times 5 - \frac12 \times 2 \times 25 = 100 - 25 = 75\ \mathrm{m}
\]

摩托车位移：
\[
x_B = v_B t = 10 \times 5 = 50\ \mathrm{m}
\]

比较：$x_A - x_B = 25\ \mathrm{m} < d = 60\ \mathrm{m}$，所以不会相撞，最小距离为 $60 - 25 = 35\ \mathrm{m}$。

\textbf{方法总结：}
\begin{enumerate}
  \item 本题中相对速度由正变负，因此速度相等时两车的接近量最大、距离最小。
  \item 对“先接近、后远离”的追及模型，可比较最大接近量与初始距离；一般问题还要检查停车、相遇和区间端点。
\end{enumerate}

\end{examplebox}

\begin{examplebox}{追及问题——图像法}

\textbf{题意分析：} 物体 $A$ 以 $v_A = 4\ \mathrm{m/s}$ 匀速运动，物体 $B$ 在 $A$ 后方 $s_0 = 8\ \mathrm{m}$ 处从静止开始以 $a = 2\ \mathrm{m/s^2}$ 匀加速追赶 $A$。求：
\begin{enumerate}
  \item $B$ 追上 $A$ 所需时间。
  \item 追上之前 $AB$ 间的最大距离。
\end{enumerate}

\textbf{物理模型：} 匀加速追匀速，$v_B < v_A$ 时距离增大，$v_B = v_A$ 时距离最大，之后 $B$ 开始缩短差距。

\textbf{解题过程：}

\begin{enumerate}
  \item 追上时 $x_B = x_A + s_0$：
  \[
  \frac12 a t^2 = v_A t + s_0,\qquad
  \frac12 \times 2 \times t^2 = 4t + 8
  \]
  \[
  t^2 - 4t - 8 = 0,\qquad
  t = \frac{4 \pm \sqrt{16 + 32}}{2} = \frac{4 \pm \sqrt{48}}{2}
  \]
  $t = 2 + 2\sqrt{3} \approx 5.46\ \mathrm{s}$（舍去负根）。

  \item 速度相等时距离最大：$v_B = a t = 4$，$t = 2\ \mathrm{s}$。

  $A$ 的位移：$x_A = 4 \times 2 = 8\ \mathrm{m}$

  $B$ 的位移：$x_B = \frac12 \times 2 \times 2^2 = 4\ \mathrm{m}$

  最大距离：$s_{\max} = s_0 + x_A - x_B = 8 + 8 - 4 = 12\ \mathrm{m}$
\end{enumerate}

在 $v$-$t$ 图中，$t_1$ 之前 $B$ 比 $A$ 慢，红色有向面积对应两者距离继续增大的量；$t_1$ 之后 $B$ 比 $A$ 快，蓝色有向面积对应 $B$ 追近的量。追上时，蓝色面积比红色面积多出的部分正好等于初始距离 $s_0$。

\begin{center}
% TikZ 图：追及问题的 v-t 图像
% Day 04：运动图像与追及相遇
% 对应例题：A 匀速 vA=4，B 从静止以 a=2 匀加速。
% t1=2 s 时速度相等，t0=2+2sqrt(3) 时 B 追上 A。
\begin{tikzpicture}[scale=0.82,>=Latex]
  \draw[->] (-0.3,0) -- (6.2,0) node[right] {$t$};
  \draw[->] (0,-0.3) -- (0,4.9) node[above] {$v$};
  \node[below left] at (0,0) {$O$};

  \coordinate (TOne) at (2,1.6);
  \coordinate (TZeroA) at (5.46,1.6);
  \coordinate (TZeroB) at (5.46,4.37);

  % 有向面积：前段 A 多走，后段 B 多追。
  \fill[red!15] (0,0) -- (0,1.6) -- (TOne) -- cycle;
  \fill[blue!15] (TOne) -- (TZeroA) -- (TZeroB) -- cycle;

  % A 匀速，B 匀加速
  \draw[very thick,red] (0,1.6) -- (5.9,1.6) node[right] {$A$};
  \draw[very thick,blue] (0,0) -- (5.9,4.72) node[above right] {$B$};

  % 速度相等与追上时刻
  \draw[dashed] (2,0) -- (TOne);
  \node[below] at (2,0) {$t_1$};
  \node[anchor=south west,fill=white,inner sep=1pt] at (2.08,1.72) {$v_A=v_B$};
  \draw[dashed] (5.46,0) -- (TZeroB);
  \node[below] at (5.46,0) {$t_0$};

  \node[red!70!black,fill=white,inner sep=1pt] at (0.82,0.78) {\scriptsize 红色：$A$ 多走};
  \node[blue!70!black,fill=white,inner sep=1pt] at (4.35,3.05) {\scriptsize 蓝色：$B$ 多追};
\end{tikzpicture}


\figcaption{追及问题的 $v$-$t$ 图：两图线之间的有向面积对应相对位移变化}\label{fig:chase-vt}
\end{center}

\textbf{方法总结：}
\begin{enumerate}
  \item 速度相等是追及问题的重要候选节点；只有结合相对速度变号和实际运动区间，才能判断它是否对应距离极值。
  \item $v$-$t$ 图中两图线围成的有向面积差对应相对位移变化。
\end{enumerate}

\end{examplebox}

\section*{七、变式训练}

\begin{enumerate}
  \item 甲车以 $v_{\text{甲}} = 15\ \mathrm{m/s}$ 匀速行驶，发现前方 $L = 40\ \mathrm{m}$ 处乙车以 $v_{\text{乙}} = 10\ \mathrm{m/s}$ 匀速行驶，甲立即以减速度大小 $a = 1\ \mathrm{m/s^2}$ 匀减速。通过计算判断是否相撞。

  \item 在平直公路上，$A$ 车在前以 $v_A = 8\ \mathrm{m/s}$ 匀速行驶，$B$ 车在后以 $v_B = 20\ \mathrm{m/s}$ 匀速行驶。当两车相距 $d = 100\ \mathrm{m}$ 时，$B$ 开始以减速度大小 $a = 2\ \mathrm{m/s^2}$ 匀减速。求：
  \begin{center}
  \begin{tikzpicture}[scale=0.8,>=Stealth]
  \draw[thick] (-0.4,0) -- (6.2,0);
  \draw[fill=gray!10] (0,0) rectangle (0.9,0.42);
  \draw[fill=gray!10] (4.0,0) rectangle (4.9,0.42);
  \node[above] at (0.45,0.42) {$B$};
  \node[above] at (4.45,0.42) {$A$};
  \draw[->,blue,thick] (0.95,0.3) -- (1.8,0.3) node[midway,above] {$v_B$};
  \draw[->,blue,thick] (4.95,0.3) -- (5.65,0.3) node[midway,above] {$v_A$};
  \draw[->,red,thick] (0.95,0.12) -- (0.35,0.12) node[midway,below] {$a$};
  \draw[<->] (0.9,-0.35) -- (4.0,-0.35) node[midway,below] {$d$};
  \node at (2.45,0.82) {同向追及};
\end{tikzpicture}

  \end{center}
  \begin{enumerate}
    \item $B$ 车能否追上 $A$？
    \item 若能，何时追上？若不能，最近距离是多少？
  \end{enumerate}

  \item 一质点沿 $x$ 轴运动，其 $v$-$t$ 图线为：$0-2\ \mathrm{s}$ 内匀速 $v = 4\ \mathrm{m/s}$，$2-4\ \mathrm{s}$ 内匀减速 $v = 4 - 2(t-2)$（从 $4$ 减到 $0$），$4-6\ \mathrm{s}$ 内反向匀加速 $v = -2(t-4)$（从 $0$ 反向加速到 $-4$）。求全程的位移和路程。
  \begin{center}
  \begin{tikzpicture}[scale=0.78,>=Stealth,line cap=round,line join=round]
  \draw[->] (0,-2.05) -- (0,2.35) node[above] {$v/(\mathrm{m\,s^{-1}})$};
  \draw[->] (0,0) -- (5.65,0) node[right] {$t/\mathrm{s}$};
  \fill[blue!10] (0,0) -- (0,1.6) -- (1.7,1.6) -- (3.4,0) -- cycle;
  \fill[red!10] (3.4,0) -- (5.1,-1.6) -- (5.1,0) -- cycle;
  \draw[very thick,blue!75!black] (0,1.6) -- (1.7,1.6) -- (3.4,0) -- (5.1,-1.6);

  \foreach \x/\lab in {1.7/2,3.4/4,5.1/6} {
    \draw (\x,0.08) -- (\x,-0.08);
    \node[below] at (\x,-0.10) {$\lab$};
  }
  \draw[dashed,gray!75] (1.7,0) -- (1.7,1.6);
  \draw[dashed,gray!75] (5.1,0) -- (5.1,-1.6);
  \node[left] at (0,1.6) {$4$};
  \node[left] at (0,-1.6) {$-4$};
  \node[font=\small,blue!70!black] at (1.25,0.62) {正位移};
  \node[font=\small,red!70!black] at (4.48,-0.72) {负位移};
\end{tikzpicture}

  \end{center}
\end{enumerate}

\begin{learningbox}{变式训练参考结果}
\begin{enumerate}
  \item 甲车与乙车速度相等需要 $5\ \mathrm{s}$，在此之前甲车最多接近
  \[
  \Delta x_{\text{相,max}}
  =(15-10)\times 5-\frac12\times1\times5^2
  =12.5\ \mathrm{m}<40\ \mathrm{m}.
  \]
  因此不会相撞，最小距离为 $27.5\ \mathrm{m}$。

  \item 相对速度 $u=v_B-v_A=12-2t$，在 $t=6\ \mathrm{s}$ 时由正变负；最大接近量为
  \[
  \Delta x_{\text{相,max}}=\frac12\times6\times12=36\ \mathrm{m}<100\ \mathrm{m}.
  \]
  因此不能追上，最近距离为 $64\ \mathrm{m}$。

  \item 正向位移为 $8+4=12\ \mathrm{m}$，反向位移为 $-4\ \mathrm{m}$，所以全程位移
  $\Delta x=8\ \mathrm{m}$，路程 $s=8+4+4=16\ \mathrm{m}$。
\end{enumerate}
\end{learningbox}

\section*{八、竞赛入门补充}

\begin{competitionbox}
\textbf{相对运动、分段制动与安全距离}

竞赛题常把追及问题转化为相对运动。取前车为参考对象，令“后车相对前车多走的距离”为 $\Delta x_{\text{相}}$，相对速度为
\[
u=v_{\text{后}}-v_{\text{前}}
\]
在每一段匀速或匀变速过程中，$u$-$t$ 图与时间轴围成的有向面积，就是该段后车相对前车多走的距离；再把各段面积按正负相加。$u>0$ 时距离缩短，$u<0$ 时距离增大；面积累计到初始距离 $d$ 时恰好追上。若 $u$ 从正变负，正面积的最大累计值给出最大追近量，从而可求最近距离。

\textbf{标准处理顺序：}
\begin{enumerate}
  \item 先按反应、制动、停车或变速时刻把过程分段；每一段写 $u$ 的表达式或画 $u$-$t$ 图。
  \item 对每一段计算有向面积，累计为相对位移，不把不同阶段的加速度混在一条公式里。
  \item 令累计相对位移与初始距离比较：小于 $d$ 未追上，等于 $d$ 临界相遇，大于 $d$ 则在该阶段内发生追上。
\end{enumerate}

\textbf{例题：} 前车以 $12\ \mathrm{m/s}$ 匀速行驶，后车以 $24\ \mathrm{m/s}$ 同向行驶。司机反应 $0.5\ \mathrm{s}$ 后才开始刹车，刹车加速度大小为 $4\ \mathrm{m/s^2}$。求刚好不追尾所需的最小初始距离。

\textbf{解：} 看相对速度 $u=v_{\text{后}}-v_{\text{前}}$。反应时间内 $u=12\ \mathrm{m/s}$，多走
\[
s_1=12\times 0.5=6\ \mathrm{m}
\]
开始刹车后，后车速度从 $24$ 降到 $12\ \mathrm{m/s}$ 需要
\[
t=\frac{24-12}{4}=3\ \mathrm{s}
\]
这段时间 $u$ 从 $12$ 均匀减到 $0$，相对速度图像是三角形，后车又多走
\[
s_2=\frac12\times 3\times 12=18\ \mathrm{m}
\]
因此最小安全距离
\[
s_{\min}=s_1+s_2=24\ \mathrm{m}
\]

\textbf{竞赛提示：} 有反应时间、分段刹车或两车都变速时，先画 $u$-$t$ 图。它只关心“距离怎样变”，常比直接列两辆车的位置方程更清楚；但最后仍要把“面积是否达到初始距离”翻译回实际的相遇或安全距离结论。

\medskip
\textbf{进一步：两车同时制动但停车时刻不同}

前车初速度为 $18\ \mathrm{m/s}$，以 $6\ \mathrm{m/s^2}$ 的加速度大小制动；后车初速度为 $24\ \mathrm{m/s}$，同时以 $4\ \mathrm{m/s^2}$ 的加速度大小制动。两车均沿同一直线同向运动，求避免追尾所需的最小初始距离。

前车在 $3\ \mathrm{s}$ 时先停下，后车在 $6\ \mathrm{s}$ 时停下，因此必须分成两段：
\[
u(t)=v_{\text{后}}-v_{\text{前}}
=
\begin{cases}
(24-4t)-(18-6t)=6+2t, & 0\le t\le 3\ \mathrm{s},\\
24-4t, & 3<t\le 6\ \mathrm{s}.
\end{cases}
\]

\pagebreak[4]
\begin{center}
% TikZ 图：两车同时制动时的分段相对速度 u-t 图
\begin{tikzpicture}[scale=0.8,>=Latex,line cap=round,line join=round]
  \draw[->] (-0.25,0) -- (6.85,0) node[right] {$t/\mathrm{s}$};
  \draw[->] (0,-0.2) -- (0,3.55) node[above] {$u/(\mathrm{m\,s^{-1}})$};
  \node[below left] at (0,0) {$O$};

  \fill[orange!18] (0,0) -- (0,1.5) -- (3,3) -- (3,0) -- cycle;
  \fill[blue!14] (3,0) -- (3,3) -- (6,0) -- cycle;
  \draw[very thick,orange!85!black] (0,1.5) -- (3,3);
  \draw[very thick,blue!75!black] (3,3) -- (6,0);

  \draw[dashed,gray!75] (3,0) -- (3,3);
  \draw[dashed,gray!75] (0,1.5) -- (0.35,1.5);
  \draw[dashed,gray!75] (0,3) -- (3,3);
  \node[left] at (0,1.5) {$6$};
  \node[left] at (0,3) {$12$};
  \node[below] at (3,0) {$3$};
  \node[below] at (6,0) {$6$};
  \node[above] at (3,3) {前车停车};
  \node[anchor=north] at (6,-0.42) {后车停车};
  \node[orange!80!black] at (1.55,0.78) {$\Delta x_1$};
  \node[blue!70!black] at (4.15,0.82) {$\Delta x_2$};
\end{tikzpicture}


\figcaption{两车同时制动的 $u$-$t$ 图：前车先停车后，相对速度表达式随之改变}\label{fig:relative-braking-ut}
\end{center}

图中两块面积分别为
\[
\Delta x_1=\frac{6+12}{2}\times3=27\ \mathrm{m},\qquad
\Delta x_2=\frac12\times3\times12=18\ \mathrm{m}.
\]
后车停止前相对速度始终为正，因此两车一直在接近，所需最小初始距离为
\[
d_{\min}=\Delta x_1+\Delta x_2=45\ \mathrm{m}.
\]

\textbf{模型提升：} 一旦某辆车停车，它的速度就保持为零，后续不能继续沿用原来的匀减速表达式。分段点不仅来自“开始刹车”，也来自任一物体停车、改变加速度或改变运动方向。
\end{competitionbox}

\section*{九、易错点}

\begin{mistakebox}
\begin{enumerate}
  \item \textbf{混淆 $x$-$t$ 图和 $v$-$t$ 图：} 读图第一步先圈出纵轴物理量和单位。纵轴是位置 $x$ 时，斜率才是速度；纵轴是速度 $v$ 时，斜率才是加速度，有向面积才是位移 $\Delta x$。

  \item \textbf{误以为 $v$-$t$ 图交点就是相遇：} 先分清“速度相等”和“位置相同”。$v$-$t$ 图交点只表示速度相等，判断相遇还要比较位移和初始距离。

  \item \textbf{忽略面积的正负：} 先按 $t$ 轴上方为正、下方为负求有向面积，得到的是位移；若题目问路程，再把各段面积的绝对值相加。

  \item \textbf{追及问题忘记初始距离：} 列式前先画一条位置轴，标出前车、后车和初始距离 $d$。再写 $x_{\text{后}}=x_{\text{前}}+d$ 或等价关系，不要默认两车从同一点出发。

  \item \textbf{把速度相等一律当作距离极值：} 速度相等只是区间内部的候选条件，还要检查相对速度是否变号，并与初始、停车、相遇和考察终点比较。

  \item \textbf{匀减速追及忘记停车时间：} 分别算各车停车时间，再判断停车前是否追上；物体停下后速度保持为零，不能把匀减速公式继续延长到负速度。
\end{enumerate}
\end{mistakebox}

\section*{十、今日小结}

\begin{enumerate}
  \item $x$-$t$ 图的斜率 $=$ 速度，$v$-$t$ 图的斜率 $=$ 加速度，有向面积对应区间位移 $\Delta x$。
  \item 两种图像的“水平线”含义不同：$x$-$t$ 水平 $\to$ 静止，$v$-$t$ 水平 $\to$ 匀速。
  \item 追及问题的核心：位置相等表示相遇；速度相等时相对速度为零，是距离极值的候选条件。
  \item 判断是否追尾：分段累计最大接近量，并与初始距离比较；同时检查停车、相遇和各段端点。
  \item 图像法和代数法互相校验：先用图像判断趋势、关键时刻和面积，再用方程计算数值。
\end{enumerate}

\textbf{一分钟复盘：} 任选一条 $x$-$t$ 直线和一条 $v$-$t$ 直线，在图旁各写三项：斜率代表什么、面积是否有物理意义、交点表示什么。能把这三项分清，图像法的主线就清楚了。
